% Publications - a page listing research articles written using content in the thesis

\pdfbookmark[1]{Publications}{Publications} % Bookmark name visible in a PDF viewer

\chapter*{Publications} % Publications page text
\addcontentsline{toc}{chapter}{\textsc{Publications}}

This thesis encompasses and extends results from my previously published works. Here, I describe how these publications relate to the thesis as a whole. The detailed structure of the thesis is summarized in \autoref{chapter:introduction}. The groundwork for several concepts examined here was laid in my Ph.D. thesis proposal~\cite{Smijakova2021}.

My first publication~\spacedlowsmallcaps{[CMSB'20]} is based on the results that I have achieved in my master thesis. Since all subsequent work is done on formal methods applied to systems biology, this publication gave me a solid foundation in the field. Moreover, it introduced me to the topic of colored model checking based on BDDs, which is related to the methods I have developed in this thesis. The context of these methods is described in \autoref{chapter:software}.

Then I focused on the topic of control of partially specified Boolean networks, which is the backbone of this thesis. The two published works~\spacedlowsmallcaps{[Mathem'21, Biosys'23]} focus on source-target control, and their results are demonstrated in \autoref{chapter:source_target_control}.

The findings from these works inspired me to investigate phenotype control of partially specified Boolean networks because it alleviates the limitations of source-target control in the partially specified setting. The results achieved with this method were published in~\spacedlowsmallcaps{[CMSB'23]} and here are described in \autoref{chapter:phenotype_control_psbn}.

The methods I developed in~\spacedlowsmallcaps{[Biosys'23, CMSB'23]} were then included in the AEON.py tool. This tool was presented as the Application note~\spacedlowsmallcaps{[Bioinf'23]}. Later, the control methods were incorporated into the frontend of the AEON tool, as is presented in~\spacedlowsmallcaps{[CMSB'25]}. The relevant parts of the tools are described here in \autoref{chapter:software}.

Finally, I proposed a novel framework to design control experiments for refinement of the partially specified Boolean networks.
This framework and relevant case studies are presented in a journal article, that is accepted but not yet published at the time of the submission of this thesis~\spacedlowsmallcaps{[TOAPPEAR]}. The framework along with other case studies for control of partially specified Boolean networks, is summarized in \autoref{chapter:applications}.

The following list chronologically enumerates the publications I co-authored. For each publication, I describe my involvement in their preparation, and an approximate percentage indicating my contribution.


% The crazy hspace stuff makes causes the titles to break lines when we need them to.

\medskip
\noindent\PrimaryColor{\spacedlowsmallcaps{[TOAPPEAR]}}\hspace{3pt} \emph{Control-Guided Refinement of Partially Specified Bool\-e\-an Networks: Applications to RTK Signalling} (E. Šmijáková, L. Brim, S. Pastva, D. Šafránek)


\begin{quotation}
	\emph{An original journal paper. At the time of this thesis submission, the paper was accepted and is going to appear in the Bioinformatics journal. I have conceptualized the novel framework for control-guided refinement of partially specified Boolean networks and methods for controlling partially specified networks to oscillatory phenotypes. I conducted experiments to demonstrate the methods as well as wrote the major part of the paper. (70\%)}
\end{quotation}


\medskip
\noindent\PrimaryColor{\spacedlowsmallcaps{[CMSB'25]}}\hspace{3pt} \emph{AEON 2025: Robust Control of Partially-Specified Boolean Networks} (V. Veselý, E. Šmijáková, S. Pastva, N. Beneš, D. Šafránek) \cite{Vesely2025}


\begin{quotation}
	\emph{A tool conference paper for Computational Methods in Systems Biology conference. The main author is a student whom I supervised. The student developed an extension of the AEON tool frontend based on the methods I have previously developed. I have provided consultations to the student, contributed to the implementation and testing of the tool, and written a part of the paper. (30\%)}
\end{quotation}


\medskip
\noindent\PrimaryColor{\spacedlowsmallcaps{[CMSB'23]}}\hspace{3pt} \emph{Phenotype Control of Partially Specified Boolean Net\-works} (L. Brim, S. Pastva, D. Šafránek, E. Šmijáková)~\cite{Benes2023b}


\begin{quotation}
	\emph{A regular conference paper for Computational Methods in Systems Biology conference. I have conceptualized the novel type of control problem with better usability than the previous ones. I have proposed and implemented the methods. I designed and conducted case studies case studies that demonstrated the applicability of the developed methods. I contributed substantially to the writing. (70\%)}
\end{quotation}


\medskip
\noindent\PrimaryColor{\spacedlowsmallcaps{[Biosys'23]}}\hspace{3pt} \emph{Temporary and Permanent Control of Partially Specified Boolean Networks} (L. Brim, S. Pastva, D. Šafránek, E. Šmijáková) \cite{brim2022}

\begin{quotation}
	\emph{A full-length journal article in Biosystems journal. I proposed, developed, and implemented the novel methods introduced in the paper. I designed and executed the experiments demonstrating the capabilities of the methods and did an extensive part of the writing. (70\%)}
\end{quotation}


\medskip
\noindent\PrimaryColor{\spacedlowsmallcaps{[Bioinf'23]}}\hspace{3pt} \emph{AEON.py: Python library for attractor analysis in asynchronous Boolean networks
} (N. Beneš, L. Brim, O. Huvar, S. Pastva, D. Šafránek, E. Šmijáková)~\cite{aeonpy}

\begin{quotation}
	\emph{An application note published in the Bioinformatics journal. I was the author of the control methods included in the published tool paper. I have described the technical details of these methods. (30\%)}
\end{quotation}

\medskip
\noindent\PrimaryColor{\spacedlowsmallcaps{[Mathem'20]}}\hspace{3pt} \emph{Parallel One-Step Control of Parametrised  Networks} (L. Brim, S. Pastva, D. Šafránek, E. Šmijáková)~\cite{Brim2021}

\begin{quotation}
	\emph{An article published in a special issue on Boolean Networks Models in Science and Engineering in Mathematics journal. I have conceptualized the novel problem of the control of partially specified Boolean networks. I have proposed the algorithm that solves this problem and collaborated on its implementation and wrote a major part of the man\-u\-script. (70\%)}
\end{quotation}


\medskip
\noindent\PrimaryColor{\spacedlowsmallcaps{[CMSB'20]}}\hspace{3pt} \emph{Parallel Parameter Synthesis for Multi-affine Hybrid Systems from Hybrid CTL Specifications} (E. Šmijáková, S. Pastva, D. Šaf\-rá\-nek, L. Brim)~\cite{Smijakova2020}

\begin{quotation}
	\emph{A regular conference paper for Computational Methods in Systems Biology conference. The article presents the results achieved in my master's thesis~\cite{Smijakova2019}. I have proposed and implemented a method based on a previous work of the SyBiLa research group. I significantly participated in the manuscript writing. (60\%)}
\end{quotation}

% \hspace{4pt}


% \endgroup

%Some ideas and figures have appeared previously in the following publications:\\

%\noindent Put your publications from the thesis here. The packages \texttt{multibib} or \texttt{bibtopic} etc. can be used to handle multiple different bibliographies in your document.

%\begin{refsection}[ownpubs]
%    \small
%    \nocite{*} % is local to to the enclosing refsection
%    \printbibliography[heading=none]
%\end{refsection}

%\emph{Attention}: This requires a separate run of \texttt{bibtex} for your \texttt{refsection}, \eg, \texttt{ClassicThesis1-blx} for this file. You might also use \texttt{biber} as the backend for \texttt{biblatex}. See also \url{http://tex.stackexchange.com/questions/128196/problem-with-refsection}.